\sloppy
\section{Background and Related Work}
\label{sec:background}

\subsection{Test Generation}

Unit test generation can be broadly categorized into traditional approaches and Deep Learning (DL)-based approaches. 
Traditional approaches, exemplified by tools like Evosuite~\cite{fraser_evosuite_2011} and Randoop~\cite{pacheco_randoop_2007}, utilize a series of software analysis techniques, including search-based~\cite{fraser_evosuite_2011}, random-based~\cite{pacheco_randoop_2007}, model checking~\cite{enoiu_automated_2016, gargantini_using_1999}, and symbolic execution~\cite{pasareanu_combining_2008,xie_symstra_2005}, to generate unit tests that achieve high code coverage and mutation scores.

DL-based approaches, by using pre-trained models, treat the test generation task as a neural machine translation problem, where the input primarily consists of the focal methods and the output is the generated unit tests. 
Tufano et al.~\cite{tufano_unit_2021} introduce an approach named \textsc{AthenaTest} to generate unit tests, which use a BART architecture model for a two-step training procedure on a large processed dataset. Alagarsamy et al.~\cite{alagarsamy_a3test_2023} propose an approach named A3Test. In addition to using a PLBART model, A3Test implements simple post-processing steps such as verifying naming consistency, correcting incomplete parentheses, and refining test signatures to increase the correctness rate.
More recently, with the significant impact of closed-source LLMs like ChatGPT, many researchers have explored using a prompt engineering approach to generate unit tests. Yuan et al.~\cite{yuan_no_2024} propose ChatTester, which combines an initial test generator with an iterative test refiner to produce unit tests. Chen et al.~\cite{chen_chatunitest_2024} introduce ChatUniTest, which generates unit tests under the Generation-Validation-Repair framework. 
Gu et al.\cite{gu2024testart} introduce TestART, a framework that generates high-quality unit test cases by integrating traditional template-based program repair\cite{zhang2023gamma} with the generative capabilities of LLMs within a \textit{generation-and-repair} mechanism.
Zhang et al.~\cite{zhang2024testbench} introduce TestBench, the first benchmark tailored to class-level LLM-based test case generation, involving 108 Java programs from 9 large-scale GitHub projects, three distinct prompt types and five key evaluation aspects of test cases.
Besides, some empirical studies~\cite{yang_empirical_2024, schaefer_empirical_2024} have emerged, comparing the effectiveness of different LLMs with various prompt strategies.







\subsection{Assertion Generation}


Assertion generation is primarily approached using two approaches~\cite{he2024empirical}: DL-based approaches and Information Retrieval (IR)-based approaches.
DL-based approaches treat assertion generation as a Neural Machine Translation (NMT) task. 
Watson et al.~\cite{watson_learning_2020} pioneer the use of DL-based approaches in assertion generation with their proposal of $ATLAS$. 
Later, IR-based approaches are applied to the assertion generation task as well, which use retrievers to find similar assertions based on focal-test information and then correct the retrieved assertions by replacing incorrect tokens with the correct ones based on context.
Yu et al.~\cite{yu_automated_2022} propose an IR-based method, introducing IR-based assertion retrieval ($IR_{ar}$) and retrieved-assertion adaptation ($RA_{adapt}$) approaches. $IR_{ar}$ retrieves the most similar assertion to the given focal-test, while $RA_{adapt}$ adjusts tokens based on context to refine the assertion. Yu et al. further propose an integrated approach (abbreviated as $Integration$), which comprehensively considers assertions generated by both IR-based and DL-based approaches.
Recently, hybrid approaches~\cite{sun_revisiting_2023} combining DL-based methods and IR-based methods have been proposed and proven effective. 
These combined approaches leverage the strengths of both approaches, resulting in more meaningful and useful assertions for developers.


\subsection{Test Evolution}

Hu et al.~\cite{hu_identify_2023} first use a DL-based method to generate updated test cases for the test evolution task in \textsc{Ceprot}. \textsc{Ceprot} focuses on the changes (i.e., edit sequences) between two versions of the production code and attempts to transfer these changes to the outdated test code to generate updated test code. Yaraghi et al.~\cite{yaraghi_automated_2024} propose an approach named \textsc{TaRGet}, which employs a two-step process to fine-tune a model based on the dataset. \textsc{TaRGet} focuses on repairing outdated test code. The input to the model is not limited to the related production code but also considers potential context from all changes in the project between the two versions. Recently, Liu et al.~\cite{liu_augmenting_2024} propose \textsc{Synbciatr}, which focuses on syntactic breaking changes and attempts to automatically repair outdated test cases via precise and concise Test-Repair-Oriented Contexts (TROCtx) construction.




\subsection{Large Language Models}

LLMs are large-scale models pre-trained on massive textual corpora~\cite{ouyang_training_2022, raffel_exploring_2020, vaswani_attention_2017}, demonstrating strong capabilities across various Natural Language Processing (NLP) tasks~\cite{zhao_survey_2024} and Software Engineering (SE)~\cite{zhang2023survey_se,wang_software_2024}. 
LLMs are primarily built on the Transformer~\cite{vaswani_attention_2017} architecture, which includes an encoder for input representation and a decoder for output generation.
Based on the structure, LLMs can be categorized into three types: (1) encoder-only models (e.g., CodeBERT~\cite{feng_codebert_2020}), designed for understanding tasks; (2) encoder-decoder models (e.g., CodeT5~\cite{wang_codet5_2021}), designed for translation tasks; and (3) decoder-only models (e.g., CodeGen~\cite{nijkamp_codegen_2023}), designed for generation tasks.

To improve the performance of LLMs on unseen downstream tasks, researchers often fine-tune pre-trained models using task-specific datasets~\cite{zhang2023pre,yuan2022circle}.
However, with recent advancements in LLM capabilities, particularly in in-context learning, there is a growing trend towards using prompt engineering as an alternative approach to handling downstream tasks~\cite{wang_software_2024}.
While commercial LLMs continue to dominate the top of the leaderboards recently, an increasing number of open-source models, such as CodeLlama~\cite{roziere_code_2024} and DeepSeek-Coder~\cite{guo_deepseek-coder_2024}, are emerging and demonstrating strong performance across various tasks.
